\section{The unital complexes through height three}
\label{low:section}

We first construct the lower complexes, retaining a binary two-sided unit.
The algebraic input is the May upper bound of Lemma~\ref{may:page}:
a bidegree containing no initial-page monomial has zero cobar cohomology.
The two vanishing families below serve different parts of the induction:
the first makes each $v_n$ survive to a self-map, and the second permits
the two unit maps to extend to a binary product over every pair of cells.

\begin{lemma}[Lower obstruction groups]
\label{low:algebraic-vanishing}
The following cobar cohomology groups vanish.
\begin{enumerate}
 \item For $1\leq n\leq3$ and $k\geq1$, the group with coefficients
 $R/I_n$ in bidegree
 \[
  (s,w)=(qk+1,u_n+k).
 \]
 \item For $0\leq n\leq3$, choose $e_i\in\{0,1,2\}$ for $0\leq i\leq n$
 and put
 \[
  U=\sum_{i=0}^n e_i u_i,\qquad
  \rho=\sum_{i=0}^n e_i\geq2.
 \]
 For every $k\geq1$, the group with coefficients $R/I_{n+1}$ in bidegree
 \[
  (s,w)=(qk+1-\rho,U+k)
 \]
 is zero.
\end{enumerate}
\end{lemma}

\begin{proof}
Include the first case in the notation of the second by setting
$U=u_n$ and $\rho=0$, while retaining its indicated coefficient comodule.
If the normalized cochain group in degree $(s,w)$ is nonzero,
Lemma~\ref{fw:normalized-bound} gives
\begin{equation}
 2015k\leq U+\rho-1.
 \label{low:normalized-cutoff}
\end{equation}
The exact values
\[
 (u_0,u_1,u_2,u_3,u_4)=(0,1,1010,1{,}019{,}091,1{,}028{,}262{,}820)
\]
give Table~\ref{low:cutoffs}.  Each upper bound is obtained by substituting
the displayed maxima in Equation~\eqref{low:normalized-cutoff} and taking
the integer part.  A zero upper bound on $k$ excludes all $k\geq1$.

\begin{table}[htbp]
\centering
\begin{tabular}{lrrrr}
\toprule
Case & Maximum $U$ & Maximum $\rho$ & Maximum $k$ & Maximum $w$\\
\midrule
Product, $n=0$ & $0$ & $2$ & $0$ & ---\\
Product, $n=1$ & $2$ & $4$ & $0$ & ---\\
Product, $n=2$ & $2022$ & $6$ & $1$ & $2023$\\
Product, $n=3$ & $2{,}040{,}204$ & $8$ & $1012$ & $2{,}041{,}216$\\
$v_1$ target & $1$ & $0$ & $0$ & ---\\
$v_2$ target & $1010$ & $0$ & $0$ & ---\\
$v_3$ target & $1{,}019{,}091$ & $0$ & $505$ & $1{,}019{,}596$\\
\bottomrule
\end{tabular}
\caption{Preliminary bounds for all lower obstruction bidegrees.}
\label{low:cutoffs}
\end{table}

These bounds also justify the additive comparison in the adjacent cochain
degrees.  For either remaining product case, $w<u_{n+1}$, so no positive
coefficient or coordinate $t_i$ with $i\geq n+1$ can occur.  The retained
coproducts are additive modulo $I_{n+1}$, and the last assertion of
Lemma~\ref{fw:range-comparison} applies in every degree.  For a remaining
$v_n$ target, a positive coefficient costs at least $u_n$, leaving at most
$k$ weight for $d\geq s-1=qk>k$ reduced bars when
$d\in\{s-1,s,s+1\}$.  Such a coefficient cannot occur.  A bar containing
$t_n$ or a higher coordinate leaves at most $k$ for the other
$d-1\geq qk-1>k$ bars and cannot occur either.  Hence the same three
cochain degrees are additive.  In all cases, both adjacent differentials
are identified, so the desired cohomology is identified with that for $P$.

Every remaining weight in Table~\ref{low:cutoffs} is at most $2{,}041{,}216$,
whereas
\[
 p^3=1{,}027{,}243{,}729.
\]
A May letter covering position $3$ or above has weight at least $p^3$.
Thus only positions $0,1,2$ can occur. We use the common exterior and
polynomial alphabet of Equation~\eqref{may:low-alphabet}, with the interval
coverage conventions of Section~\ref{may:preliminaries}.
Allowing this entire alphabet only enlarges the upper bound if fewer
coordinates were retained in the preceding comparison.

Suppose a monomial of the desired bidegree exists.  Let $m\leq6$ be its
number of exterior letters and $L=(s-m)/2$ its total polynomial exponent.
Every polynomial letter has weight at least $p$, so
\[
 U+k\geq pL=p(p-1)k+\frac{p(1-\rho-m)}2.
\]
For all profiles under consideration,
$U\leq2p^2+4p+6$ and $\rho\leq8$.  Consequently
\begin{equation}
 (p^2-p-1)k
 \leq U+\frac{p(\rho+5)}2
 \leq 2p^2+\frac{21}2p+6.
 \label{low:polynomial-cutoff}
\end{equation}
If $k\geq3$, the left side exceeds the last expression by at least
\[
 p^2-\frac{27}2p-9=\frac{2{,}008{,}901}{2}>0.
\]
It follows that $1\leq k\leq2$.

Write $U=F_0+pF_1+p^2F_2$.  In a product profile the digits are
\[
 F_0=e_1+e_2+e_3,\qquad F_1=e_2+e_3,\qquad F_2=e_3,
\]
with missing $e_i$ set to zero.  For a $v_n$ target they are initial
strings of ones.  In every case,
\begin{equation}
 6\geq F_0\geq F_1\geq F_2\geq0,
 \qquad F_1\leq4,\quad F_2\leq2.
 \label{low:digits}
\end{equation}
Let $H_r$ count the exterior letters covering position $r$, and let
$n_B,n_C$ be the polynomial exponents of $B,C$.  Adding weights gives
the exact identity
\begin{equation}
 F_0+k+pF_1+p^2F_2
  =H_0+p(H_1+L-n_C)+p^2(H_2+n_B+n_C).
 \label{low:weight-equation}
\end{equation}
We next extract its digits without an unproved no-carry assumption.

First, $0\leq H_0\leq3$ and $1\leq F_0+k\leq8<p$.  Reduction modulo $p$
therefore gives $H_0=F_0+k$.  Also $B,C$ each cost at least $p^2$, so
\[
 p^2(n_B+n_C)\leq w\leq2p^2+4p+8<3p^2;
\]
the last strict inequality has difference
$p^2-4p-8=1{,}014{,}037>0$.  Hence $n_B+n_C\leq2$.

Put $\delta=(m+\rho-1)/2$.  It is an integer because $s-m$ is even.
For a target with $\rho=0$, that parity forces $m$ odd, hence $m\geq1$;
for a product profile $\rho\geq2$.  Thus $\delta\geq0$.  The bounds
$m\leq6$, $\rho\leq8$ and the same parity give $\delta\leq6$.
Now $L=(p-1)k-\delta$.  Subtracting the position-zero equation from
Equation~\eqref{low:weight-equation}, dividing by $p$, and rearranging gives
\[
 H_1-k-\delta-n_C-F_1
   =p(F_2-k-H_2-n_B-n_C).
\]
Since $H_1\leq4$, the left side lies in $[-14,3]$.  Its only possible
multiple of $p=1009$ is zero.  We have therefore proved the three exact
relations
\begin{equation}
 H_0=F_0+k,\qquad
 H_2+n_B+n_C=F_2-k,\qquad
 H_1=F_1+k+\delta+n_C.
 \label{low:coverage-relations}
\end{equation}

Among the exterior letters covering position zero, only $x,a$ fail to
cover position two.  Thus $H_2\geq H_0-2$.  Equations~\eqref{low:digits}
and~\eqref{low:coverage-relations} imply
\[
 F_0-F_2+2k\leq2.
\]
As $k\geq1$, equality is forced: $k=1$ and
$F_0=F_1=F_2=f$.  All the intervening inequalities are equalities, giving
\[
 n_B=n_C=0,\qquad H_0=f+1,\qquad H_2=f-1.
\]
The bounds on $H_0,H_2$ leave only $f=1,2$.

If $f=1$, the exterior letters $x,a$ must occur and $c,b,z$ must be
absent.  Only $y$ is optional, so $H_1\leq2$ and $m\geq2$.  On the other
hand, the last relation in Equation~\eqref{low:coverage-relations} gives
$H_1=2+\delta>2$, since
$\delta=(m+\rho-1)/2>0$.  If $f=2$, the letters $x,a,c$ must occur and
$b,z$ must be absent.  Again only $y$ is optional, so $H_1\leq3$ and
$m\geq3$, whereas the same relation gives $H_1=3+\delta>3$.
Both possibilities are impossible.  There is no monomial on the May
upper-bound page, and hence no cobar cohomology in any of the stated
bidegrees.
\end{proof}

\begin{lemma}[Finite local cell attachments]
\label{low:finite-localization}
A spectrum obtained from finitely many $p$-local sphere cells is the
$p$-localization of a finite CW spectrum.
\end{lemma}

\begin{proof}
Proceed one attachment at a time.  If the previously constructed spectrum
is the localization of a finite spectrum $Z$, an attaching map belongs to
$\pi_d(Z_{(p)})=\pi_d(Z)\otimes\mathbb Z_{(p)}$ and can be written $a/m$
with $p\nmid m$.  After localization, multiplication by $m$ on its sphere
source is an equivalence.  The localized cone of the integral attaching
map $a$ is therefore equivalent to the cone of $a/m$.  Induction over
the finitely many cells proves the assertion.
\end{proof}

\begin{proposition}[Construction through $V(3)$]
\label{low:V3}
For $0\leq n\leq3$ there is a finite $p$-local cell spectrum $V(n)$ with
a bottom-cell unit $\eta_n:S_{(p)}\to V(n)$ and a binary multiplication
\[
 \mu_n:V(n)\wedge V(n)\longrightarrow V(n)
\]
whose two unit identities hold.  Its cells are indexed by subsets of
$\{0,\ldots,n\}$, in dimensions
\begin{equation}
 \sum_{i\in S}(qu_i+1),
 \label{low:cell-dimensions}
\end{equation}
and its $BP$-homology is the canonical quotient comodule $R/I_{n+1}$,
generated in degree zero by the bottom unit.  In particular, we may fix
\begin{equation}
 Y=V(3),\qquad BP_*Y=R/I_4,
 \label{low:Y}
\end{equation}
with such a binary two-sided unital multiplication.
\end{proposition}

\begin{proof}
Start with $V(0)$, the cofiber of $p:S_{(p)}\to S_{(p)}$.  Since
multiplication by $p$ is injective on $R$, its $BP$-homology is $R/(p)$,
with the canonical quotient coaction and bottom generator.  Its cells
are in dimensions zero and one, as in
Equation~\eqref{low:cell-dimensions}.

We give first the product-extension argument, which applies to this
initial complex and to every later complex once its asserted homology
and cells have been obtained.  The subcomplex
\[
 U_n=(S_{(p)}\wedge V(n))
       \mathop{\cup}_{S_{(p)}\wedge S_{(p)}}
       (V(n)\wedge S_{(p)})
       \ \subseteq\ V(n)\wedge V(n)
\]
has a map to $V(n)$ given by the two identity maps, which agree on the
common bottom sphere.  Every remaining cell corresponds to a pair of
nonempty subsets of $\{0,\ldots,n\}$.  If $e_i$ counts occurrences of $i$
in that pair, its dimension is $qU+\rho$, with precisely the parameters
in the second part of Lemma~\ref{low:algebraic-vanishing}.  The obstruction
to extending across the cell lies in
$\pi_{qU+\rho-1}V(n)$.

An Adams--Novikov term in that stem has
$qw-s=qU+\rho-1$, or
\[
 w=U+k,\qquad s=qk+1-\rho.
\]
Since $2\leq\rho\leq8<q$, nonnegative filtration requires $k\geq1$.
Every corresponding $E_2$ group is zero by
Lemma~\ref{low:algebraic-vanishing}.  Strong convergence and finite
filtration, from Proposition~\ref{fw:anss}, give
$\pi_{qU+\rho-1}V(n)=0$.  Extending cell by cell therefore produces
$\mu_n$, agreeing with both identity maps on $U_n$.  This constructs the
initial product on $V(0)$ as well.

Suppose now that $V(n-1)$, including its product, has been constructed,
where $1\leq n\leq3$.  By Lemma~\ref{fw:coordinates}, $v_n$ is a primitive
element of $BP_*V(n-1)=R/I_n$.  It defines a filtration-zero $E_2$ class.
Its possible positive differential targets are exactly
Equation~\eqref{fw:primitive-targets}; the first part of
Lemma~\ref{low:algebraic-vanishing} makes them all zero.  There is no
incoming differential to filtration zero.  The edge assertion in
Proposition~\ref{fw:anss} consequently realizes $v_n$ by a homotopy map
\[
 \widetilde v_n:S^{qu_n}_{(p)}\longrightarrow V(n-1).
\]
Use the product to form
\[
 f_n=\mu_{n-1}(\widetilde v_n\wedge\operatorname{id}):
       \Sigma^{qu_n}V(n-1)\longrightarrow V(n-1),
\]
and define $V(n)$ as its cofiber.

The induced map on $BP$-homology is multiplication by $v_n$: the
coefficient pairing is $R$-bilinear, the homology is cyclic on the unit,
and $\widetilde v_n$ has the specified Hurewicz image.  This multiplication is
injective on the polynomial quotient $R/I_n$.  The cofiber exact sequence
therefore gives precisely $R/I_{n+1}$, with no additional shifted kernel
summand.  The bottom sphere map surjects onto this cyclic homology, so its
kernel $I_{n+1}$ identifies the coaction with the canonical quotient
coaction, as well as fixing the generator in degree zero.

The mapping cone adds the old cells shifted by $qu_n+1$, giving exactly
Equation~\eqref{low:cell-dimensions}.  The product-extension argument
already proved then constructs $\mu_n$.  This completes the induction.
All attachments are finite $p$-local cell attachments, and
Lemma~\ref{low:finite-localization} gives finiteness in the sense of the
stated realization problem.  No associativity of any $\mu_n$ has been
asserted or used.
\end{proof}

\begin{corollary}[The action of $p$ on $Y$]
\label{low:p-null}
For the spectrum $Y$ of Proposition~\ref{low:V3}, $\pi_0Y=\mathbb F_p$ and
$p\operatorname{id}_Y$ is nullhomotopic.
\end{corollary}

\begin{proof}
The bottom two cells form the mod-$p$ Moore spectrum.  All other cells have
dimension at least $qu_1+1=q+1>1$, so they do not change $\pi_0$.
Thus $p\eta_3=0$ in $\pi_0Y$.  Applying the left unit identity of its
binary multiplication gives
\[
 p\operatorname{id}_Y
  =\mu_3\bigl((p\eta_3)\wedge\operatorname{id}_Y\bigr)=0.
\]
\end{proof}
